% Generated draft: requires mathematical and source review.
\documentclass[11pt]{article}
\usepackage[margin=25mm]{geometry}
\usepackage{fontspec}
\setmainfont{Latin Modern Roman}
\usepackage{amsmath,amssymb,mathtools,unicode-math}
\setmathfont{Latin Modern Math}
\usepackage{enumitem,needspace,xurl,hyperref}
\setlength{\parindent}{0pt}
\setlength{\parskip}{5pt}
\setlength{\emergencystretch}{3em}
\newcommand{\N}{\mathbb{N}}
\newcommand{\Z}{\mathbb{Z}}
\newcommand{\Q}{\mathbb{Q}}
\newcommand{\R}{\mathbb{R}}
\newcommand{\C}{\mathbb{C}}
\newcommand{\F}{\mathbb{F}}
\newcommand{\A}{\mathbb{A}}
\newcommand{\PP}{\mathbb{P}}
\newcommand{\E}{\mathbb{E}}
\newcommand{\eps}{\varepsilon}
\newcommand{\dd}{\,\mathrm{d}}
\begin{document}
\section*{UnsolvedMath ID: 20003143}
\textbf{Source problem code:} catalogue entry
\subsection{Definitions and mathematical statement}
The source's complete request is ``Study collapse to a ray.'' It does not prescribe dimension, curvature normalization, completeness, pointed versus compact convergence, or a classification target. In the nonnegative-sectional-curvature workshop context, the aid supplies an explicit rotational family rather than a general classification.

Fix $n\geq2$. On $\R^n$ in polar coordinates, consider
\[
 g_i=dr^2+f_i(r)^2g_{S^{n-1}},\qquad r\geq0,
\]
where $g_{S^{n-1}}$ is the unit round metric, $f_i(r)>0$ for $r>0$, and the pole is smooth: $f_i(0)=0$, $f_i'(0)=1$, with the usual smooth odd extension at zero. Suppose
\[
 \sup_{r\geq0}f_i(r)\leq a_i,\qquad a_i\longrightarrow0.
\]
The aid reports that the radius map $\rho_i:(\R^n,g_i)\to[0,\infty)$ is an exact submetry, sending every metric ball onto the corresponding ray ball. On the ambient closed ball around the pole $o$, with its restricted ambient distance, it reports the uniform estimate
\[
 d_{\mathrm{GH}}\bigl(\overline B_{g_i}(o,R),[0,R]\bigr)
 \leq\frac{\pi a_i}{2}\qquad\text{for every }R>0.
\]
The bound is independent of $R$, and yields pointed Gromov--Hausdorff convergence to the ray based at zero. The assumption $n\geq2$ ensures connected spherical angular fibers; the same argument is not asserted for $\R^1$.

If in addition $f_i''\leq0$ and $0\leq f_i'\leq1$, the reported nonnegative sectional curvature follows from
\[
 K_{\mathrm{rad}}=-\frac{f_i''}{f_i},\qquad
 K_{\mathrm{tan}}=\frac{1-(f_i')^2}{f_i^2},
\]
with the tangential formula applying to tangential two-planes when $n\geq3$. Curvature extends smoothly across the pole under the pole conditions.

The explicit example is
\[
 f_\varepsilon(r)=\varepsilon\tanh(r/\varepsilon),\qquad\varepsilon>0.
\]
It gives complete smooth nonnegatively curved metrics on $\R^n$, $\sup f_\varepsilon=\varepsilon$, and the reported volume estimate
\[
 \operatorname{Vol}_{g_\varepsilon}\overline B(o,R)
 \leq\operatorname{Vol}(S^{n-1})\,R\varepsilon^{n-1}.
\]
Thus every fixed-radius ball collapses in volume for $\varepsilon\to0$. The fibers over $r>0$ are spheres and the fiber over zero is a point, identified as the point soul in this example. This endpoint degeneration is not a globally regular sphere bundle over the closed ray.

These hypotheses impose no common upper sectional-curvature bound. In particular, the example's smooth-pole curvature scales like $\varepsilon^{-2}$, so it must not be presented as a normalized two-sided-bounded-curvature collapse. The record gives an existence model and quantitative radius-map control, but no classification of all nonnegatively curved collapses to a ray or every possible endpoint fiber.
\subsection{Short English statement}
The aid gives complete rotational examples collapsing to a ray, with an error bound uniform over ball radii and an explicit volume rate. The endpoint fiber becomes a point. The examples do not have a uniform upper curvature bound and are not a classification of ray collapse.
\subsection{Sources}
\begin{itemize}
\item Local UnsolvedMath registry, permanent ID 20003143, source code catalogue entry; \texttt{aim-topology-notes.json}, notes[230].
\item M. Kerin's compilation \emph{Open Problems in Non-negative Sectional Curvature}, from the September 2007 AIM workshop \emph{Manifolds with Non-negative Sectional Curvature}, Problem 18: \url{https://aimath.org/WWN/nnsectcurvature/nnsectcurvature.pdf}.
\item The bounded-warping criterion and explicit hyperbolic-tangent model are supplied by the research summary without a separate authoritative external citation.
\item Machine-generated research attempt 1: partial result, open at the research run, rigor labelled ``proved,'' low novelty confidence; registry partially solved. Not peer reviewed; independently verify before citation. Registry updated 14 August 2026.
\end{itemize}
\end{document}
